\section{Experiments}
\label{sec:experiments}

\subsection{Models and Evaluation Protocol}

\paragraph{Evaluated Models.}
Across distinct generation tasks—namely Text-to-Video (T2V; Director3D~\cite{wu2024direct3dscalableimageto3dgeneration}, OpenSoraPlan~\cite{lin2024opensoraplan}, T2V-Turbo~\cite{li2024t2v-turbo}, HunyuanVideo~\cite{kong2024HunyuanVideo}), Image-to-Video (IT2V; Matrix Game2.0~\cite{he2025matrix}, Wan2.1~\cite{wan2025}, Wan2.2~\cite{wan2025}, CogVideo~\cite{hong2022CogVideo}, OpenSora~\cite{zheng2024opensora}, Cosmos~\cite{agarwal2025Cosmos}, LargeVideoPlanner~\cite{chen2025large}), and camera-controlled generation (HunyuanWorld~\cite{hyworld2025}, HunyuanGameCraft~\cite{li2025hunyuan}, ViewCrafter~\cite{yu2024viewcrafter}, Gen3C~\cite{ren2025gen3c}, Lingbot~\cite{team2026advancing}, FantasyWorld~\cite{dai2025fantasyworldgeometryconsistentworldmodeling}, WonderWorld~\cite{yu2025wonderworld})--we evaluate a total of 18 representative world models encompassing diffusion-based, autoregressive, and hybrid paradigms.

\paragraph{Evaluation Protocol.}
We comprehensively evaluate the generative capabilities of world models using our proposed benchmark, \textbf{Omni-WorldBench}. The evaluation protocol is driven by \textbf{Omni-Metric} (defined in Sec.~\ref{sec:metrics}), which encompasses 15 metrics across three distinct dimensions: (1) generated video quality, (2) interaction effect fidelity, and (3) camera and object controllability. To compute these metrics, we introduce a custom test set, \textbf{Omni-WorldSuite}. Specifically, we evaluate T2V and IT2V models using 410 diverse prompts from this suite, while camera-conditioned models are assessed using a dedicated set of 120 prompts equipped with explicit camera trajectories.


\begin{sidewaystable}[p]
\centering
\caption{Quantitative evaluation results of various models on the proposed benchmark. The metrics are grouped into Interaction Effect Fidelity, Generated Video Quality, Camera-Object Controllability, and the overall AgenticScore. The best results within each group are highlighted in bold. Avg.=average.}
\label{tab:main_results}
\sisetup{
  table-number-alignment=center,
  table-format=3.2,
  detect-weight=true,
  detect-inline-weight=math
}
\small
\setlength{\tabcolsep}{1.5pt}
\renewcommand{\arraystretch}{1.5}
\begin{adjustbox}{width=\textwidth}
\begin{tabular}{
l
*{5}{S}
*{6}{S}
c
*{3}{S}
S
}
\toprule
& \multicolumn{5}{c}{Interaction Effect Fidelity}
& \multicolumn{6}{c}{Generated Video Quality}
& \multicolumn{4}{c}{Camera-Object Controllability}
& \multicolumn{1}{c}{AgenticScore} \\
[2pt]
\cmidrule(lr){2-6}\cmidrule(lr){7-12}\cmidrule(lr){13-16}\cmidrule(lr){17-17}
Model
& {InterStab-L}
& {InterStab-N}
& {InterCov}
& {InterOrder}
& {Avg.}
& {\makecell[c]{Imaging\\Quality}}
& {\makecell[c]{Temporal\\Flickering}}
& {\makecell[c]{Content\\Alignment}}
& {\makecell[c]{Motion\\Smoothness}}
& {\makecell[c]{Dynamic\\Degree}}
& {Avg.}
& {\makecell[c]{Camera\\Control}}
& {\makecell[c]{Transitions\\Detect}}
& {\makecell[c]{Object\\Control}}
& {Avg.}
& {(\%$\uparrow$)} \\
\midrule


\rowcolor{gray!12}\multicolumn{17}{c}{\textbf{T2V}}\\
Director3D      & 73.49 & 89.24 & 44.48 & 38.41 & 61.41 & 48.90 & \textbf{99.48} & 89.87 & \textbf{99.68} & 47.31 & 77.05 & \textemdash & \textbf{99.75} & 72.07 & 85.91 & 71.00 \\
OpenSoraPlan    & 68.78 & \textbf{95.76} & 40.29 & 36.59 & 60.36 & 53.55 & 98.67 & 82.16 & 99.22 & 16.83 & 70.09 & \textemdash & 98.29 & 70.65 & 84.47 & 68.10 \\
T2V-Turbo       & \textbf{82.98} & 66.18 & 43.83 & 36.70 & 57.42 & \textbf{63.48} & 97.64 & \textbf{90.24} & 98.54 & \textbf{47.80} & \textbf{79.54} & \textemdash & 99.02 & 73.75 & 86.39 & 69.85 \\
HunyuanVideo    & 77.35 & 82.37 & \textbf{53.02} & \textbf{46.78} & \textbf{64.88} & 61.60 & 98.67 & 81.18 & 99.31 & 44.88 & 77.13 & \textemdash & 98.54 & \textbf{85.30} & \textbf{91.92} & \textbf{73.96} \\

\rowcolor{gray!12}\multicolumn{17}{c}{\textbf{IT2V}}\\
Matrix Game2.0       & 47.38 & 19.96 & 55.27 & 48.41 & 42.76 & 58.85 & 95.37 & 62.44 & 98.17 & \textbf{99.02} & 82.77 & \textemdash & 53.41 & 87.85 & 70.63 & 60.33 \\
Wan2.1             & 70.98 & 58.53 & \textbf{64.52} & \textbf{54.19} & 62.06 & 65.89 & 96.75 & 81.56 & 98.04 & 70.98 & 82.64 & \textemdash & 81.95 & \textbf{91.93} & 86.94 & 73.21 \\
Wan2.2             & 79.68 & 79.98 & 56.99 & 52.70 & \textbf{67.34} & \textbf{66.83} & 98.36 & 79.67 & 99.09 & 46.83 & 78.16 & \textemdash & 96.83 & 91.18 & 94.01 & \textbf{75.92} \\
CogVideo           & 79.03 & 79.51 & 54.80 & 48.98 & 65.58 & 61.47 & 98.04 & 79.19 & 98.84 & 29.02 & 73.31 & \textemdash & 97.56 & 87.33 & 92.45 & 73.27 \\
OpenSora           & 66.68 & 69.90 & 62.54 & 48.17 & 61.82 & 57.40 & 98.29 & \textbf{86.09} & 99.09 & 79.76 & \textbf{84.13} & \textemdash & 95.12 & 90.51 & 92.82 & 74.71 \\
Cosmos             & 79.55 & 79.63 & 53.89 & 51.81 & 66.22 & 66.30 & 98.29 & 80.93 & 99.17 & 44.15 & 77.77 & \textemdash & \textbf{98.78} & 91.01 & \textbf{94.90} & 75.42 \\
LargeVideoPlanner  & \textbf{82.15} & \textbf{87.43} & 42.84 & 45.15 & 64.39 & 66.60 & \textbf{98.99} & 77.67 & \textbf{99.36} & 32.44 & 75.01 & \textemdash & 97.32 & 89.84 & 93.58 & 73.42 \\

\rowcolor{gray!12}\multicolumn{17}{c}{\textbf{With Camera}}\\
HunyuanWorld        & 77.49 & \textbf{67.92} & 55.31 & 48.15 & \textbf{62.22} & 64.27 & \textbf{97.65} & \textbf{76.10} & \textbf{99.14} & 67.32 & 80.90 & 55.40 & 96.10 & 87.52 & 79.67 & \textbf{74.36} \\
HunyuanGameCraft  & 64.78 & 51.28 & 46.74 & 37.50 & 50.08 & 67.09 & 96.12 & 47.29 & 98.67 & 91.67 & 80.17 & 27.96 & 95.00 & 84.55 & 69.17 & 67.39 \\
ViewCrafter         & 81.15 & 4.22 & 43.19 & 41.11 & 42.42 & 61.37 & 91.01 & 49.03 & 95.40 & \textbf{100.00} & 79.36 & 42.91 & 95.00 & 86.17 & 74.69 & 65.88 \\
Gen3c               & 75.90 & 38.40 & \textbf{57.50} & \textbf{53.75} & 56.39 & 58.55 & 95.78 & 63.84 & 98.86 & 98.33 & 83.07 & 48.07 & 85.83 & 84.55 & 72.82 & 71.61 \\
Lingbot             & 74.84 & 66.59 & 45.28 & 35.28 & 55.50 & \textbf{67.65} & 96.93 & 52.83 & 98.67 & 45.83 & 72.38 & 33.97 & \textbf{98.33} & 89.76 & 74.02& 67.16 \\
FantasyWorld        & 72.66 & 55.34 & 48.40 & 41.94 & 54.59 & 64.87 & 96.32 & 56.32 & 98.68 & 73.33 & 77.90 & 42.29 & 93.33 & \textbf{90.45} & 75.36 & 69.49 \\
WonderWorld         & \textbf{84.96} & 24.89 & 51.26 & 43.84 & 51.24 & 60.40 & 92.26 & 74.22 & 99.02 & \textbf{100.00} & \textbf{85.18} & \textbf{96.12} & 73.95 & 87.33 & \textbf{85.80} & 74.02 \\
\bottomrule
\end{tabular}
\end{adjustbox}
\end{sidewaystable}




\subsection{Implementation Details}

All inference experiments are conducted using NVIDIA H20 GPUs. To ensure optimal performance and fair comparison, the software environments—specifically the Python and PyTorch versions—are strictly configured according to the official guidelines provided by each model's respective codebase. 

\paragraph{Text-to-Video (T2V) Models.}
For the T2V generation paradigm, models are conditioned solely on text prompts. Specifically, HunyuanVideo generates 91 frames at a $1280\times 720$ resolution using 50 inference steps at 16 FPS. OpenSoraPlan (v1.0.0) employs the T5-v1.1-XXL text encoder, producing 65 frames at $512\times 512$ resolution with 250 sampling steps and a classifier-free guidance (CFG) scale of 7.5 at 24 FPS. T2V-Turbo (v2-no-MG V) generates 40 frames at 8 FPS, utilizing 32 inference steps and a CFG scale of 7.5. Notably, Director3D relies on its self-predicted camera trajectories for novel view rendering, outputting $960\times 960$ resolution videos.

\paragraph{Image-to-Video (IT2V) Models.}
IT2V models utilize both a starting frame and text prompts as conditioning inputs. Both Wan2.1 (14B-720P) and Wan2.2 (A14B) generate 81 frames (5 seconds) at $1280\times 720$ resolution operating at 16 FPS; however, Wan2.1 uses 50 steps with a 5.0 guidance scale, whereas Wan2.2 uses 40 steps with a 3.5 guidance scale. Cosmos (Cosmos-predict-14B) operates at the same resolution and frame rate but outputs 77 frames using 35 steps and a guidance scale of 7. CogVideo (CogVideoX-5b-I2V) generates 49 frames at $720\times 480$ (8 FPS, 50 steps, CFG scale 6). OpenSora (v2) is configured to a 256px (16:9) resolution, yielding 129 frames at 24 FPS with 50 steps and a 7.5 CFG scale. LargeVideoPlanner leverages a base model for $832\times 480$ resolution (81 frames, 16 FPS, 40 steps) with customized history and language guidance scales of 1.5 and 2.5, respectively. Finally, Matrix-Game2.0 (universal mode) outputs $650\times 352$ videos at 16 FPS, relying on randomly generated camera trajectories.

\paragraph{Camera-Conditioned Models.}
To evaluate camera controllability, these models require explicit camera parameters or trajectories. HunyuanWorld (v1.5, Autoregressive-480P-I2V) generates videos at $800\times 496$ resolution and 16 FPS. Hunyuan-GameCraft utilizes complete pose information to generate 132 frames at $704\times 1216$ (24 FPS). ViewCrafter adopts an equidistant camera pose sampling strategy, producing 25 frames at $576\times 1024$ (8 FPS). Both Gen3C and Lingbot operate at $720\times 1280$ resolution, with their outputs consistently truncated to the first 121 frames. Furthermore, FantasyWorld ($832\times 480$) and WonderWorld ($512\times 512$) employ a frame-subsampling strategy, compressing the original 132-frame camera trajectories down to 81 frames. It is worth noting that for WonderWorld, large-scale camera motions in the dataset may occasionally result in blank frames during rendering due to incomplete point cloud coverage.


\begin{figure}
    \centering
    \includegraphics[width=\linewidth]{Images/visual-1-0323.pdf}
    \caption{\textbf{Non-camera-controlled Interaction Comparison.} Qualitative comparison of generated videos from different models under the same prompt and first-frame condition. Representative frames illustrate differences in interaction effect fidelity, motion dynamics, and scene coherence during the throwing action.}
    \label{fig:visual}
\end{figure}



\begin{figure}
    \centering
    \includegraphics[width=\linewidth]{Images/visual-2.pdf}
    \caption{\textbf{Camera-Controlled Interaction Comparison.} Qualitative Comparison of Generated Videos from Different Models under the Same Prompt, First-Frame, and Camera Trajectory Condition.}
    \label{fig:visual2}
\end{figure}



\subsection{Quantitative Evaluation Results and Analysis}
\label{subsec:Omni-Metric_exp}

This section presents a comprehensive automatic evaluation of various advanced video generation models on the proposed benchmark. As shown in Tab.~\ref{tab:main_results}, different model categories exhibit clear trade-offs among interaction fidelity, video quality, and controllability.

\textbf{Overall Performance:} Overall, the Image-to-Video (IT2V) paradigm, which incorporates richer conditional inputs like images, demonstrates the highest performance potential on the current benchmark. Wan2.2 achieves the highest overall AgenticScore across all models at 75.92\%, closely followed by Cosmos (75.42\%). Among pure Text-to-Video (T2V) models, HunyuanVideo performs the best, reaching 73.96\%. In the group supporting explicit camera control (With Camera), hunyuanworld (74.36\%) and wonderworld (74.02\%) take the lead.

\textbf{Interaction Effect Fidelity:} This dimension evaluates the stability of models in handling complex physical and logical interactions. The IT2V group shows high consistency, with Wan2.2 achieving the highest average score of 67.34\%. Notably, some models in the ``With Camera'' group exhibit a significant trade-off across different interaction sub-metrics. For instance, wonderworld scores an impressive 84.96\% on InterStab-L but drops sharply to 24.89\% on InterStab-N. This indicates that maintaining consistent underlying interaction logic while introducing complex camera scheduling remains a challenge for current models.

\textbf{Generated Video Quality:} In terms of basic visual quality, the vast majority of evaluated models have reached extremely high levels in Temporal Flickering and Motion Smoothness (mostly exceeding 95.00\%). However, there is a substantial variance in the Dynamic Degree across models, which constitutes a core differentiator in generation capabilities. ViewCrafter and WonderWorld achieve a perfect score of 100.00\%, while other models in the same group vary significantly. Therefore, the major differences across models no longer mainly come from temporal smoothness, but rather from content alignment and dynamic responsiveness.

This metric directly reflects the models' ability to precisely control specific elements. Camera-aware methods show clear advantages here. WonderWorld demonstrates an overwhelming advantage with an explicit Camera Control score of 96.12\%, far surpassing other models in the same category. Meanwhile, HunyuanWorld obtains the best average controllability score of 79.67\% in its group. Furthermore, regarding Object Control, Cosmos (94.90\%) and Wan2.2 (94.01\%) excel in the IT2V group.

\textbf{Summary:} Current models are already strong in conventional video quality metrics, but still show clear limitations in action-conditioned world evolution, causal interaction consistency, and joint camera-object control. These results highlight the importance of evaluating world models beyond passive video quality and toward agent-centric interactive generation.






\subsection{Qualitative Evaluation}

\paragraph{Visual Comparison of T2V and IT2V Models.} To provide a concrete illustration of our evaluation on interaction effect fidelity and motion dynamics, we present a qualitative comparison in Fig.~\ref{fig:visual}. The models are evaluated under a challenging \textbf{Level-2 interaction prompt} that requires generating a baseball player performing a powerful throw. As shown in the visual sequences, Wan2.2~\cite{wan2025} demonstrates superior performance in this scenario; it successfully synthesizes a complete, anatomically reasonable pitching motion while maintaining the athlete's structural integrity and scene coherence throughout the video. In stark contrast, Matrix-Game2.0~\cite{he2025matrix} struggles significantly with this complex physical interaction. The generated action is not only incomplete but also suffers from severe temporal degradation, culminating in the catastrophic collapse and complete disappearance of the human figure in the final frames. These qualitative observations—particularly the stark disparities in physical interaction and temporal consistency—are highly consistent with the quantitative results presented in Sec.~\ref{subsec:Omni-Metric_exp}, further validating the effectiveness of our Omni-Metric evaluation framework.




\paragraph{Visual Comparison of Camera-Conditioned Models.} In our qualitative analysis, we categorize this example as a \textbf{Level-1 interaction} (camera view trajectory control: left strafe). As shown in Fig.~\ref{fig:visual2}, HunyuanWorld~\cite{hyworld2025} exhibits relatively stable performance throughout the sequence. In contrast, ViewCrafter~\cite{yu2024viewcrafter} introduces a spurious building that appears out of nowhere, degrading visual consistency and leading to a lower score. This qualitative observation is consistent with our quantitative evaluation results presented in Sec.~\ref{subsec:Omni-Metric_exp}, further validating the effectiveness of our Omni-Metric evaluation framework.




























% \subsection{Automatic Evaluation on VBench}
% \label{subsec:auto_vbench_worldscore}

% Tab. \ref{tab:vbench} reports the automatic evaluation results on VBench. 
% Overall, world-model–based methods (TIC2V) achieve higher overall scores compared with pure T2V and I2V approaches, indicating stronger performance on dynamic reasoning and scene evolution. 
% For instance, HYGC1.0 achieves the best overall score (80.46), followed by Wan2.1 (78.51) and Lingbot-World (77.98). 
% In terms of metric breakdown, most models obtain very high scores on \textit{Subject Consistency}, \textit{Background Consistency}, and \textit{Motion Smoothness} (generally above 95\%), suggesting that current video generators can reliably maintain visual continuity. 
% However, the scores on \textit{Dynamic Degree} show larger variance across models, where world-model–based approaches consistently outperform T2V methods, reflecting their stronger capability in modeling complex motion and interaction dynamics. 
% In contrast, T2V models remain competitive on appearance-related metrics such as \textit{Aesthetic Quality} and \textit{Imaging Quality}, highlighting their advantage in visual fidelity rather than dynamic reasoning. 
% We provide WorldScore~\cite{duan2025worldscore} results in App.~\ref{app:worldscore}.


% \begin{table}[t]
% \centering
% \caption{Automatic evaluation on VBench (\%).}
% \label{tab:vbench}

% \sisetup{
%   table-number-alignment=center,
%   table-format=2.2,
%   detect-weight=true,
%   detect-inline-weight=math
% }
% \setlength{\tabcolsep}{3pt}
% \renewcommand{\arraystretch}{1.1}

% \begin{adjustbox}{width=\linewidth}
% \begin{tabular}{l *{7}{S}}
% \toprule
% Model
% & \multicolumn{1}{c}{\makecell{Subject\\Consistency$\uparrow$}}
% & \multicolumn{1}{c}{\makecell{Background\\Consistency$\uparrow$}}
% & \multicolumn{1}{c}{\makecell{Motion\\Smoothness$\uparrow$}}
% & \multicolumn{1}{c}{\makecell{Dynamic\\Degree$\uparrow$}}
% & \multicolumn{1}{c}{\makecell{Aesthetic\\Quality$\uparrow$}}
% & \multicolumn{1}{c}{\makecell{Imaging\\Quality$\uparrow$}}
% & \multicolumn{1}{c}{\makecell{Overall $\uparrow$}} \\
% \midrule

% \rowcolor{gray!12}\multicolumn{8}{l}{\textbf{T2V}}\\
% Cogvideo      & 95.31 & 96.69 & 98.89 & 17.10 & 49.98 & 59.72 & 69.62 \\
% % Opensora      & 91.61 & 95.09 & 99.01 & 71.94 & 46.98 & 58.99 & 77.27 \\
% OpenSoraPlan  & 95.84 & 97.38 & 99.37 & 10.97 & 51.51 & 56.03 & 68.52 \\
% HunyuanVideo  & 95.19 & 96.61 & 99.51 & 26.45 & 47.34 & 60.24 & 70.89 \\
% % Wan2.2        & 61.15 & 79.50 & 91.60 & 99.03 & 37.08 & 49.14 & 69.58 \\
% T2v-turbo     & 97.39 & 97.39 & 98.65 & 37.74 & 55.84 & 65.05 & 75.34 \\
% Director3d    & 92.28 & 95.58 & 99.68 & 40.00 & 51.49 & 50.30 & 71.56 \\
% Wan2.1        & 92.38 & 95.03 & 98.34 & 68.39 & 52.48 & 64.46 & 78.51 \\

% \rowcolor{gray!12}\multicolumn{8}{l}{\textbf{I2V}}\\
% Cosmos        & 93.93 & 95.75 & 99.30 & 33.55 & 55.72 & 66.13 & 74.06 \\
% LVP           & 97.25 & 97.11 & 99.44 & 30.32 & 54.80 & 65.68 & 74.10 \\
% % MG2.0 & 79.22 & 88.58 & 98.26 & 99.35 & 45.09 & 58.63 & 78.19 \\

% \rowcolor{gray!12}\multicolumn{8}{l}{\textbf{TIC2V}}\\
% GEN3C                 & 87.89 & 92.60 & 98.88 & 80.77 & 46.73 & 57.06 & 77.32 \\
% HYGC1.0 & 87.50 & 91.58 & 98.80 & 94.62 & 49.41 & 60.88 & 80.46 \\
% Lingbot-World         & 87.51 & 92.38 & 98.57 & 73.08 & 52.35 & 63.97 & 77.98 \\
% ViewCrafter           & 86.24 & 91.24 & 96.48 & 83.85 & 47.35 & 57.15 & 77.05 \\
% HYWP1.5        & 91.32 & 93.90 & 99.24 & 62.58 & 52.61 & 63.00 & 77.11 \\
% \bottomrule
% \end{tabular}
% \end{adjustbox}
% \end{table}



% \subsection{Human Study}
% \label{subsec:human_study}
% Our evaluation prompts are drawn from our Omni-WorldBench, covering seven physical principles and task-oriented interaction scenarios, and organized by interaction complexity (1/2/3 interactions), resulting in 34 prompts in total (30 general and 4 task-oriented). For pairwise comparison, general prompts randomly sample 4 models from 12 non-camera models and conduct pairwise matches repeated 5 times, while task-oriented prompts sample 4 models from 13 candidates with the same procedure. This results in 1020 A/B comparisons evaluated by 14 human annotators. We measure the correlation between automatic metrics and human preferences using Spearman correlation ($\rho$). Higher values indicate stronger agreement with human judgments. A correlation above 0.5 suggests that the metric is reasonably aligned with human perception of interaction quality. Full instructions and Quality control (QC) procedures are provided in App.~\ref{app:human_study}.

% \subsection{Qualitative Evaluation}

% Tab.~\ref{tab:main_results} reports automatic evaluation results on our Omni-Metric. Overall, models with stronger conditioning signals (e.g., I2V and world-model–based methods) achieve higher AgenticScore (AS) than pure T2V models. For instance, LVP~\cite{chen2025large} obtains the best overall score (AS=50.80), followed by HYWP1.5~\cite{hyworld2025} (50.08) and Director3D~\cite{wu2024direct3dscalableimageto3dgeneration} (49.50). In terms of metric breakdown, I2V models generally perform better on \textbf{Interaction Effect Fidelity (IEF)}, especially on the stability-related metric ISL, indicating improved preservation of target entities during interaction. World-model–based methods show relatively stronger results on Spatiotemporal Causal Coherence (SCC), suggesting better modeling of action–state causal relations. In contrast, several T2V models achieve competitive performance on GeneralScore (GS) components such as motion smoothness and scene consistency, reflecting stronger visual generation quality but weaker interaction reasoning. These results highlight that different model families exhibit complementary strengths, and the proposed \textbf{Omni-Metric} provides a unified framework to reveal such differences across interaction fidelity, consistency, causal coherence, and visual quality. As shown in Fig.~\ref{fig:visual}, LVP demonstrates strong interaction performance.